\chapter{Network and security}\label{cap:security}

Security comes in three layers, each against a different risk. On a trusted
local network the first one is enough, and it costs nothing; the others allow
the same device to be exposed on a public network without changing the
protocol.

\begin{center}
\begin{tikzpicture}[node distance=4mm]
  \node[blk, minimum width=44mm, fill=m2teal!8] (t) {\textbf{Token}\\\scriptsize who are you?};
  \node[blk, minimum width=44mm, right=of t, fill=m2blue!8] (s) {\textbf{TLS}\\\scriptsize is anyone listening?};
  \node[blk, minimum width=44mm, right=of s, fill=m2amber!12] (h) {\textbf{Defences}\\\scriptsize can you be knocked over?};
  \node[lbl, below=1mm of t] {always · no cost};
  \node[lbl, below=1mm of s] {optional · $\sim$1 core};
  \node[lbl, below=1mm of h] {always · no cost};
\end{tikzpicture}
\end{center}

\section{Access token}

Every request except \ep{GET}{/v1/health} must carry the token:

\begin{lstlisting}[language=shell]
curl -H "Authorization: Bearer $(cat token)" http://192.168.18.5:8080/v1/device
\end{lstlisting}

\begin{itemize}
  \item It is a random 256-bit secret, generated at first installation, in
    \cmd{/etc/mpeg2fpgad/token} (readable only by the service and root).
  \item It is compared in constant time, so that timing the response gives
    nothing away.
  \item To rotate it:
    \cmd{ssh root@board python3 -m mpeg2fpgad -{}-rotate-token}, then restart
    the service.
\end{itemize}

\begin{advertencia}[Warning]
Without TLS the token travels unencrypted: anyone who can observe the
network's traffic can copy it. On a trusted network that is acceptable; on a
public one it is not.
\end{advertencia}

\section{TLS}

With \cmd{-{}-tls} the service serves HTTPS on port 8443 with a self-signed
certificate generated on the board. The client trusts the certificate's
\emph{fingerprint}, obtained over a secure channel, rather than a certificate
authority:

\begin{lstlisting}[language=shell]
ssh root@board python3 -m mpeg2fpgad --show-fingerprint
sha256:3f9a...
\end{lstlisting}

\begin{lstlisting}[language=Python]
dev = Device("192.168.18.5", token=token, tls_fingerprint="sha256:3f9a...")
\end{lstlisting}

If the fingerprint does not match, the library refuses to connect
(\cmd{FingerprintMismatch}): that is what stops a third party from
impersonating the device.

The cost of TLS is measured in the specifications chapter. For a public
deployment with more throughput, the alternative is to terminate TLS on a
machine in front of the device (a \emph{reverse proxy}) and keep the board on
a private network.

\section{Basic defences}

\begin{itemize}
  \item The service runs as an unprivileged user, with the file system
    read-only, and is given access only to the decoder files it uses.
  \item Ten wrong tokens in a row from one address block it for 60 seconds
    (response 429).
  \item At most 8 simultaneous connections; a connection that sends nothing
    for 30 seconds is closed.
  \item The diagnostic endpoints are off unless \cmd{-{}-debug} is given.
  \item The clip only reaches the decoder: nothing the client sends is used
    as a file name or a command. A malformed clip may decode badly, but it
    does not affect the device.
\end{itemize}
